\chapter{约束优化 MPC 全身控制与学习方法}
\label{ch:advanced}

前面的控制律多数写成一条显式公式。遇到力矩上限、关节边界、接触不滑动和多个任务同时存在时，可以把“希望做什么”和“必须满足什么”写进同一个优化问题。本章不把新名词当成自动升级，而是从能手算的二次问题出发，说明这些方法增加了什么能力以及需要承担什么代价。

\section{二次规划从一个变量开始}

\textbf{优化变量}是要求解的数；\textbf{目标函数}衡量好坏；\textbf{约束}规定可接受的范围。\textbf{二次规划（quadratic programming，QP）}常写成
\begin{equation}
\min_z\ \tfrac12z^\mathsf THz+c^\mathsf Tz,
\qquad Az=b,\quad Gz\le h.
\end{equation}
当 $H\succeq0$ 时目标是凸的，可行集合由仿射等式和线性不等式给出，所以局部最优也是全局最优；要保证解唯一，还需要在可行方向上有足够的正定性。凸不代表永远有可行解，也不代表求解时间总能满足控制周期。

先看 $\min_z\tfrac12(z-2)^2$，约束 $-1\le z\le1$。无约束驻点满足 $z-2=0$，得到 $z=2$，但不可行。可行区间内最优解为 $z=1$。这个约束在最优点成为等号，叫\textbf{活跃约束}。把输出先算成二再截成一，在这个简单标量问题中恰好得到约束最优解；在多维、有耦合和未来动态的情形中，逐分量截断通常不等价于重新优化。

\subsection{拉格朗日乘子和约束反作用}

约束为 $a^\mathsf Tz=b$ 时，引入乘子 $\lambda$：
\[
\mathcal L=\tfrac12z^\mathsf THz+c^\mathsf Tz+\lambda(a^\mathsf Tz-b).
\]
对 $z$ 和 $\lambda$ 求导，可写成
\begin{equation}
\begin{bmatrix}H&a\\a^\mathsf T&0\end{bmatrix}
\begin{bmatrix}z\\\lambda\end{bmatrix}
=\begin{bmatrix}-c\\b\end{bmatrix}.
\end{equation}
对不等式还需非负乘子与互补条件。力学中的约束反力和优化中的拉格朗日乘子有密切联系：二者都在保证约束成立时抵消不允许的变化，但具体单位和符号由约束写法决定。

\section{MPC 为什么只执行第一步}

\textbf{模型预测控制（model predictive control，MPC）}每个周期用当前测量作为起点，预测未来 $N$ 步，优化输入序列，再只执行第一步。下一次新测量到来后重新求解，因此它仍然是反馈控制。

对线性离散模型 $x_{k+1}=Ax_k+Bu_k$，可构造
\begin{align}
\min_{u_0,\ldots,u_{N-1}}\quad
&\sum_{k=0}^{N-1}(x_k^\mathsf TQx_k+u_k^\mathsf TRu_k)+x_N^\mathsf TPx_N,\\
\text{满足}\quad
&x_{k+1}=Ax_k+Bu_k,\quad x_0=x_{\rm measured},\\
&u_{\min}\le u_k\le u_{\max},\quad G_xx_k\le h_x.
\end{align}
这里可把 $x$ 定义为相对恒定参考的状态误差；移动参考时需要一致地扩展参考和动力学。$Q,R,P$ 是权重矩阵，不是系统物理惯量。

将预测递推展开：$x_1=Ax_0+Bu_0$，$x_2=A^2x_0+ABu_0+Bu_1$。堆叠以后 $X=\mathcal Ax_0+\mathcal BU$。代入二次代价，得到输入序列 $U$ 的 QP，这说明“预测”怎样变成普通矩阵计算。

\subsection{两步 MPC 的完整手算}

取一个离散位置误差模型 $e_{k+1}=e_k+u_k$，其中 $u$ 表示每拍的参考增量，单位也是弧度；它不是电机力矩。初始 $e_0=1$，两步预测，代价
\begin{equation}
J=e_1^2+e_2^2+u_0^2+u_1^2.
\end{equation}
代入 $e_1=e_0+u_0,e_2=e_0+u_0+u_1$，展开得
\begin{equation}
J=2e_0^2+4e_0u_0+2e_0u_1+3u_0^2+2u_0u_1+2u_1^2.
\end{equation}
两个偏导置零，除以二后为
\[
3u_0+u_1=-2e_0,\qquad u_0+2u_1=-e_0.
\]
因此 $u_0=-0.6e_0,u_1=-0.2e_0$，在 $e_0=1$ 时为 $(-0.6,-0.2)$，预测误差为 $(0.4,0.2)$，总代价为 $0.6$。

现在加每拍约束 $|u_k|\le0.3$。第一个输入不能取 $-0.6$。在 $u_0=-0.3$ 的边界上，对 $u_1$ 求导，得 $u_1=-(e_0+u_0)/2=-0.35$，仍超限，因此最优输入为 $(-0.3,-0.3)$。预测误差为 $(0.7,0.4)$，总代价为
\[
0.7^2+0.4^2+0.3^2+0.3^2=0.83.
\]
实际只执行 $u_0=-0.3$。若下一拍测得 $e=0.7$，重新优化后第一步仍为 $-0.3$，第二步变为 $-0.2$。若有扰动使实测变成 $0.8$，优化起点就改用 $0.8$，而不是盲目沿用原计划。

\begin{figure}[htbp]\centering
\includegraphics[width=0.94\linewidth]{figures/rt_mpc.pdf}
\caption{上述两步 MPC 的实际闭环过程和初始预测。输入限值已进入优化；每次只执行第一步，使后续计划可以随测量修正。}
\end{figure}

\subsection{什么时候 MPC 不只是截断 LQR}

一维无状态约束的例子看起来像限幅反馈。更一般时，预测器会考虑\textbf{未来刹车距离}。例如关节离上限仅 $0.1\,\mathrm{rad}$，当前速度 $1\,\mathrm{rad/s}$，最大减速度 $2\,\mathrm{rad/s^2}$，最短制动距离为 $1^2/(2\times2)=0.25\,\mathrm{rad}$，已经无法在限位前停住。正确设计应更早限制速度，不能等到位置超限才截断目标。

无约束、时不变线性二次问题可得到状态反馈形式；无限时域下与 LQR 联系起来。在有限时域、终端条件合适时，也可能恰好得到相同反馈。MPC 的价值不只在有无解析解，而在预测、约束和每拍更新的组织方式。约束问题也可以在特定低维条件下预先求成分段仿射的显式 MPC，因此“有约束就绝不可能有显式解”也过于绝对。

\section{稳定 可行与按时求解是三个要求}

有一个可行的优化解，不自动证明闭环稳定。终端代价、终端集合、递归可行性和模型误差都会影响理论结论。\textbf{递归可行性}是指这次可行并执行第一步后，下次仍能找到可行解。\textbf{软约束}引入松弛变量并惩罚违约量，能提高数值可行性，但把碰撞或硬件极限随意软化会改变任务要求。

若求解超时，上一周期计划的下一步只有在仍有效、状态误差可接受且符合约束时才能使用。切回某个 PD 也必须处理参考、积分状态和切换瞬间力矩，不能把“有兜底”当作稳定性证明。需要分别记录求解状态、迭代次数、最坏耗时、约束残差和实际跟踪偏差。

\section{WBC 把任务和物理约束写在一起}

\textbf{全身控制（whole-body control，WBC）}是一类多任务控制框架。固定基座机械臂、双臂和浮动基座人形都可以采用类似思路，但驱动和接触结构不同。常见优化变量是关节广义加速度 $\ddot q$、驱动力矩 $\tau$ 和接触力 $f$：
\begin{align}
\min\quad &\sum_i\|J_i\ddot q+\dot J_i\dot q-a_i^*\|_{W_i}^2
+\|\tau\|_{R_\tau}^2,\\
\text{满足}\quad &M\ddot q+h=S^\mathsf T\tau+J_c^\mathsf Tf,\\
&J_c\ddot q+\dot J_c\dot q=0,\\
&\tau_{\min}\le\tau\le\tau_{\max},\qquad f\in\mathcal K.
\end{align}
$h=C\dot q+g$，需要时还包括明确定义的其它已知项；$S$ 选择实际驱动的自由度；$J_c$ 是保持接触的雅可比。$\|z\|_W^2=z^\mathsf TWz$。任务加速度可由参考加速度加误差反馈得到。

在\textbf{当前构型、当前接触模式固定}时，动力学和加速度约束对这些优化变量是线性的。若把摩擦锥用线性不等式近似，就得到常见的凸 QP。若同时优化未来构型、接触时序或采用精确二阶摩擦锥，则问题类别会改变，不能再笼统称整件事为一个线性约束 QP。

\begin{examplebox}[接触力为什么会改变关节所需力矩]
在本书 2R 姿态 $q=(0,\pi/2)$，
\[
J=\begin{bmatrix}-0.5&-0.5\\0.5&0\end{bmatrix},\qquad
g=(7.3575,0)^\mathsf T.
\]
静态时环境向上托住末端，$f=(0,4)^\mathsf T\,\mathrm N$。
由 $g=\tau+J^\mathsf Tf$，
\[
J^\mathsf Tf=(2,0)^\mathsf T\,\mathrm{N\,m},\quad
\tau=(5.3575,0)^\mathsf T\,\mathrm{N\,m}.
\]
环境帮助承担了一部分重力，驱动力矩应减小。若把环境对机器人的力和机器人对环境的力混用，结果会错误地增加到 $9.3575\,\mathrm{N\,m}$。
\end{examplebox}

\section{任务权重不等于严格优先级}

同一变量 $z$，任务一希望 $z=1$，任务二希望 $z=-1$。加权目标
\begin{equation}
J=(z-1)^2+w(z+1)^2
\end{equation}
的导数为 $2(z-1)+2w(z+1)$，最优解
\begin{equation}
z^*=(1-w)/(1+w).
\end{equation}
取 $w=4$，$z^*=-0.6$，两个任务都没被精确满足。把任务一的权重加得很大，结果会接近一，但有限权重仍是一种折中。

严格分层则先确定第一任务的最优可行集合，再在这个集合中优化第二任务。在本例第一层要求 $z=1$，已经没有剩余自由度，第二层只能接受误差二。这与七轴机械臂用剩余零空间处理肘部任务的逻辑一致，只是带不等式约束时需要更一般的求解结构。

\begin{figure}[htbp]\centering
\includegraphics[width=0.87\linewidth]{figures/rt_task_weights.pdf}
\caption{两个互相冲突任务的加权最优解。权重连续改变的是折中位置，并不会突然产生严格任务优先级。}
\end{figure}

\section{浮动基座与摩擦锥}

人形机器人基座的六个运动自由度没有直接电机驱动，所以不能把固定基座机械臂的所有自由度都当作可独立施加力矩。地面反力通过接触约束影响基座运动。

对接触法向 $f_z\ge0$，各向同性库仑摩擦锥为
\begin{equation}
\sqrt{f_x^2+f_y^2}\le\mu f_z.
\end{equation}
取 $f_z=100\,\mathrm N,\mu=0.5$，切向 $(40,30)\,\mathrm N$ 的长度为 $50\,\mathrm N$，恰在边界。若只分别写 $|f_x|\le50,|f_y|\le50$，就会错误允许 $(50,50)$，其长度为 $70.71$，已超出真实摩擦锥。采用 $|f_x|+|f_y|\le\mu f_z$ 是一种保守的内接菱形近似，却会排除部分真实可行力。使用哪一种近似，需要连同保守性一起说明。

\section{学习方法放在哪一层}

\textbf{强化学习（reinforcement learning，RL）}用策略 $\pi_\theta$ 从观测产生动作，训练目标常是折扣累计奖励
\begin{equation}
J(\theta)=\mathbb E\left[\sum_{k=0}^{T}\gamma^k r_k\right],\qquad0\le\gamma<1.
\end{equation}
奖励是人为定义的训练目标，不等于物理能量或稳定性函数。三步奖励都为一、$\gamma=0.9$ 时，回报为 $1+0.9+0.81=2.71$；无限重复才趋于 $1/(1-0.9)=10$。\cite{sutton}

机器人策略可以输出关节目标、力矩残差、阻抗参数或上层动作。它们对应不同的闭环。若输出关节目标，底层 PD 有助于执行参考，但不能保证任意策略目标序列都稳定、无碰撞或不超限。若输出残差 $\tau=\tau_{\rm base}+\Delta\tau_\theta$，即便残差有幅值限幅，也可能持续注入能量，需要研究频率、方向和状态依赖。

对数据驱动辨识，神经网络或高斯过程可以拟合物理模型残差，例如
\begin{equation}
\tau=Y(q,\dot q,\ddot q)\beta+f_\theta(q,\dot q,\text{温度}).
\end{equation}
先保留明确的刚体结构，再让学习模型表达难建模部分，便于检查重力、摩擦和外推行为。但如果训练数据只覆盖某个姿态与温度，不能据此保证全部工作空间。评价至少要分开训练轨迹、验证轨迹和新的运行条件；低训练误差不是因果辨识或闭环稳定的证明。

\section{计算练习与解答}

\subsection*{题一 无约束两步 MPC}
沿用本章 $J=e_1^2+e_2^2+u_0^2+u_1^2$，若 $e_0=0.5$，求输入和预测误差。

解：$u_0=-0.6e_0=-0.3,u_1=-0.2e_0=-0.1$；$e_1=0.2,e_2=0.1$。总代价 $0.04+0.01+0.09+0.01=0.15$，与初始误差平方缩放规律一致。

\subsection*{题二 权重折中}
要求两个目标 $1$ 和 $-1$ 的加权最优位置为 $0.8$，求 $w$。

解：$0.8=(1-w)/(1+w)$，得 $0.8+0.8w=1-w$，所以 $w=1/9$。第一任务权重是第二的九倍，但第一任务仍有 $0.2$ 的误差。

\subsection*{题三 摩擦可行性}
$f_z=80\,\mathrm N,\mu=0.4$，切向力 $(20,20)\,\mathrm N$ 是否在真实摩擦锥内？是否满足本章的内接菱形近似？

解：切向长度 $\sqrt{800}=28.284\,\mathrm N$，上限 $32\,\mathrm N$，所以真实摩擦锥可行。菱形近似要求 $|20|+|20|=40\le32$，不满足，说明该线性近似更保守。

\section*{本章与原稿的对应}
本章整合 [A，第60--67页] 与 [B，第65--68、92--93页] 的优化和学习内容。补出两步 MPC 的约束手算、任务权重反例、接触力矩和平面摩擦锥计算；修正“底层 PD 自动给学习策略提供稳定保证”等不成立的泛化。
